% HD-S36. GPT-6 Astra Pro, /root/sixfolds; inherited effort not exposed.
% 21 November 2026. Exact induced block and normalizer reduction.
\subsection{HD-S36: Gaussian reflection blocks}
\label{sec:gaussian-reflection-block-scope}

Let \(G\subset GL_m(K)\), \(K=\mathbb Q(i)\), be a finite
absolutely irreducible actual matrix group with character field \(K\).
Use its natural primitive rational group-algebra Jacobian factor of
a connected full-\(G\) sphere cover, with the specified Gaussian
corner and restricted theta form. For a signature-\((3,3)\) sixfold
period family of rank nine, the following conclusions hold.

If \(G\) contains a reflection, its ordinary Hermitian determinant
is \([-1]\) or \([-3]\). This assertion has no restriction on
\(m\). In particular every such source with \(m\le8\) has one
of those two determinant classes. The latter conclusion includes
finite groups not generated by reflections and nonscalar overgroups.
The proof retains the actual block, covering, polarization, and
corner maps.

\paragraph{1. The normal root blocks and the exact induced representation.}
Let \(W\ne1\) be the normal subgroup generated by all reflections
in \(G\). With a positive \(G\)-invariant \(K\)-Hermitian form
\(h\), take the connected nonorthogonality components of their
\(K\)-root lines as in HD-S35. Their spans are mutually orthogonal,
are irreducible reflection modules, and sum to \(V\). The whole
\(G\) permutes them transitively. Choose one span \(U\), put
\(r=\dim_KU\), let \(H\) be its stabilizer, and put
\(\sigma=\rho|_U:H\to GL(U)\). This representation is absolutely
irreducible because the reflection subgroup on \(U\) is so.
Choose left coset representatives \(t_1=1,\ldots,t_a\) for
\(G/H\). The explicit map
\[
 \bigoplus_{j=1}^aU\longrightarrow V,
 \qquad (u_j)_j\longmapsto\sum_j\rho(t_j)u_j,
 \qquad m=ar,                                                \tag{S36.1}
\]
is an isomorphism of induced \(K\)-representations. If
\(gt_j=t_kh\), the corresponding block of its matrix is
\(\sigma(h)\). The form is the orthogonal sum of copies of
\(h_U=h|_U\), with no scalar on any block.

The induced trace formula shows that every character value of
\(\rho\) is a sum of character values of \(\sigma\).
Conversely \(\sigma\) has actual \(K\)-matrices. Hence
\[
       K=\mathbb Q(\operatorname{tr}\rho)
        \subseteq\mathbb Q(\operatorname{tr}\sigma)
        \subseteq K.                                       \tag{S36.2}
\]
The block therefore has the same exact Gaussian character field.
Let \(H_0=\ker\sigma\) and \(G'=H/H_0\); the resulting faithful
\(G'\)-representation is the one used below.

\paragraph{2. Actual quotient curves, parabolic pairing, and primitive factors.}
For the original cover \(\widetilde C\to\mathbb P^1\), put
\[
       Y=\widetilde C/H,\qquad C'=\widetilde C/H_0,
       \qquad C'\longrightarrow Y\ \text{with group }G'.
                                                               \tag{S36.3}
\]
The finite map \(\pi:Y\to\mathbb P^1\) realizes (S36.1) as the
finite pushforward of the block local system. In the left-deck
convention the quotient sheets are \(H\backslash G\); the sheet
\(Ht_j^{-1}\) corresponds to the block labelled \(t_jH\)
in (S36.1). More explicitly, with \(\pi_U:V\to U\) the
orthogonal projection, the literal Shapiro map of HD-S30 is
\[
 F_v(x)=\pi_U\rho(x)v,\qquad F_v(hx)=\sigma(h)F_v(x),
 \qquad F_{\rho(g)v}(x)=F_v(xg).
\]
Its inverse is
\(F\mapsto\sum_{x\in H\backslash G}\rho(x)^{-1}F(x)\);
each summand is independent of its coset representative.
The equality of this inverse with (S36.1) follows by taking
\(x=t_j^{-1}\). Thus right monodromy on the quotient sheets
and the left-coset induced matrices are identified by inversion,
including their actual continuation maps.
At a puncture, a cycle of sheets has invariant vectors precisely
when its vector is fixed by the product of the block monodromies
around that cycle. Thus the invariant stalks on both sides agree,
including ramified sheets. This identifies their middle extensions.
Finite proper pushforward has no positive higher direct image on
these sheaves, since its fibres are finite; equivalently use the
same finite-sheet cellular map and its local invariant stalks.
It gives the actual isomorphism on first parabolic cohomology.
Cup product followed by trace of sheets is the orthogonal sum in
(S36.1), so this isomorphism preserves the coefficient cup form.
The parabolic dimension of the block system is consequently six.

There is a still more direct equality of the abelian factors.
Choose any nonzero \(v\in U(K)\), write \(q=h(v,v)>0\), and
use the normalized rational matrix-coefficient idempotents
\[
\begin{split}
 e_G&=\frac{2m}{|G|q}\sum_{g\in G}
        \operatorname{Re}h(v,\rho(g^{-1})v)\,g,\\
 e_H&=\frac{2r}{|H|q}\sum_{h\in H}
        \operatorname{Re}h_U(v,\sigma(h^{-1})v)\,h.
\end{split}                                                  \tag{S36.4}
\]
If \(g\notin H\), the vector \(\rho(g^{-1})v\) is in an
orthogonal block, so its coefficient is zero. Since
\(m/|G|=r/|H|\), these are the same literal element of
\(\mathbb Q[G]\):
\[
                             e_G=e_H.                        \tag{S36.5}
\]
Schur orthogonality makes both the primitive projector on the chosen
line and its conjugate. Thus, after a common integer clears their
denominators, their image is the same abelian subvariety of
\(J\widetilde C\), with the same restricted theta form.
The two real Schur factors also agree exactly:
\[
                  \frac{2|G|}{mq}=\frac{2|H|}{rq}.            \tag{S36.6}
\]
No choice of a norm-two coefficient vector is needed here.

Let \(\alpha:\widetilde C\to C'\) be the quotient map.
The normalized idempotent \(e_H\) is the product of the normalized
\(H_0\)-average with the corresponding primitive idempotent in
\(\mathbb Q[G']\). Therefore its image is the image of the
actual pullback of that primitive factor \(B'\subset JC'\).
The identities
\[
             \alpha_*\alpha^*=[|H_0|],\qquad
             (\alpha^*)^*\theta_{J\widetilde C}
                        =|H_0|\theta_{JC'}                  \tag{S36.7}
\]
show that \(\alpha^*:B'\to B\) is a polarized isogeny.
The common \(K\)-corner is identified through (S36.2), (S36.4),
and the matrix-coefficient map. In the left-deck convention its
class-sum action is the conjugated action of HD-S35, so the
identification is \(K\)-linear. The determinant multiplier
\(|H_0|^6=N(|H_0|^3)\) is a norm.

These identities hold in families after a finite base change fixing
the finite monodromy labels and homology lattices. The period map of
\(B\) is the period map of \(B'\), composed with this fixed rational
Hodge isomorphism and the actual parameter map. Consequently its
rank is bounded by the rank and parameter dimension for the block
cover (S36.3). Additional parameters of the cover
\(\widetilde C\to C'\) do not alter that factor; the equality
of idempotents and (S36.7) prove this assertion.

\paragraph{3. The genus and branch-parameter bounds, including low genus.}
Let \(g=g(Y)\), and let \(b'\) count only the nontrivial local
monodromies of the faithful cover \(C'\to Y\). Their codimensions
\(c'_j\) satisfy
\[
                   6=r(2g-2)+\sum_{j=1}^{b'}c'_j.            \tag{S36.8}
\]
For \(r\ge2\) the global invariants and dual invariants vanish.
The fixed finite monodromy choices are discrete and finite, so the
coefficient-cover moduli is finite over the moduli of \(b'\)-pointed
curves of genus \(g\). Its dimension is \(3g-3+b'\) in the stable
cases. For \(g=1,b'>0\) it is \(b'\), and for \(g=1,b'=0\)
it is at most one. In genus zero a nontrivial irreducible coefficient
system of rank at least two has \(b'\ge3\), since the fundamental
group with at most two punctures is cyclic or trivial; its dimension
is \(b'-3\). Thus no stack automorphism dimension has been omitted.

For \(r=2\), the bounds are respectively \(7,6,5\) at
\(g=0,1,2\), and larger genus is impossible by (S36.8).
Hence this case cannot give rank nine. For \(r=3\), genus one
has dimension at most six, genus two has no branch points and
dimension three, and larger genus is impossible. In genus zero,
rank nine forces exactly twelve branch values and all \(c'_j=1\).
The full coefficient group \(G'\) is then generated by its inertia
reflections; HD-G25 (or the exact rank-three source analysis HD-S31)
applies to its natural Gaussian factor.

For every \(r\ge4\), genus at least two is impossible since
\(2r>6\), and genus one has dimension at most six. Thus only
\(g=0\) remains. When \(r=4\), the complete arbitrary finite-group
result HD-S35 applies. When \(r=1\), already (S36.1) is the
actual \(K\)-line system of HD-S34, so its split determinant
conclusion applies directly without a genus estimate.

\paragraph{4. All higher-rank reflection normalizers.}
It remains to treat \(r\ge5\). The image \(W_U\) of \(W\) on
\(U\) is a normal irreducible reflection subgroup of \(G'\).
The exact Gaussian restriction of Cohen's classification in HD-G25
leaves only
\[
 A_r,\ B_r,\ D_r,\ E_6,E_7,E_8,
             \qquad G(4,p,r),\quad p=1,2,4,                 \tag{S36.9}
\]
where the exceptional real types occur only in their indicated
ranks. The nonreal rank-five and rank-six primitive cases have the
explicit non-Gaussian traces already excluded in HD-G25; there are
no further primitive exceptions in higher rank. As before, its
finite \(K\)-intertwiner identifies the given restriction with the
standard model, and (G25.3a) retains the exact polarized isogeny.

The monomial normalizer arguments of HD-S35 work in every rank
\(r\ge3\). For \(B_r\), the two reflection classes have sizes
\(r\) and \(r(r-1)\), recovering the coordinate lines. For
\(G(4,1,r)\), use its order-four reflection lines. For
\(G(4,2,r)\), use the class sizes \(r\) and \(2r(r-1)\),
with the same explicit third-coordinate compensation in the
phase-changing diagonal conjugation. For \(G(4,4,r)\), the
\(1/\sqrt2\)-angle graph again has one four-cycle for each pair
of coordinates. Its component spans are the coordinate pair-planes,
and their nonzero pairwise intersections are precisely the coordinate
lines. Thus the whole normalizer is monomial on a literal
\(K\)-line system in these cases; HD-S34 applies.

For completeness the real normalizers needed for the other cases
are obtained from an actual chamber and diagram map. A unitary
complex normalizer of a real absolutely irreducible \(W_U\) is
\(aR\), with \(R\) real orthogonal and \(a\) a root of unity,
by the \(g^{-1}\bar g\) calculation in HD-S35. For the simply-laced
root systems here, \(R\) permutes the full normalized roots. One
way to check that no extra reflection hyperplane is missing is this:
a reflection in the Weyl group fixing a point outside the root
hyperplanes would stabilize its chamber, contradicting the simple
transitivity on chambers. Hence every reflection in that group is
a root reflection, and conjugation by \(R\) permutes their roots.

The primary root-system statements are Humphreys,
\emph{Introduction to Lie Algebras and Representation Theory},
GTM 9, Theorem 10.3(a),(b),(e), printed pp.51--52, and Section 12.2,
printed pp.65--66, including Table 1 on p.66. The latter gives
\(\operatorname{Aut}(\Phi)=W\rtimes\Gamma\), with \(\Gamma\)
the actual permutations of a fixed simple-root basis preserving its
Dynkin diagram. The original chapter text and printed pp.65--66
were reopened. A primary scan is
\url{https://math.uci.edu/~brusso/humphreys.pdf}.
The relevant map can also be seen directly: join chambers by a path
crossing one wall at a time; reflection in that wall exchanges the
adjacent chambers. Thus some \(w\in W\) carries \(R\)'s chamber
back to the fixed one, and \(wR\) permutes its inward simple-root
normals. This is the specified orthogonal diagram map.

For \(D_r\), \(r\ge5\), the only diagram symmetry interchanges
the two short fork ends. In the usual simple basis
\[
 e_1-e_2,\ldots,e_{r-1}-e_r,\ e_{r-1}+e_r,
\]
it is the literal sign change \(e_r\mapsto-e_r\). Hence its
real normalizer is contained in the signed permutation group
\(W(B_r)\), and its Gaussian finite normalizer is contained in
\(\mu_4W(B_r)\): all these real matrices are rational, so the
finite phase lies in \(K\cap\mu_\infty=\mu_4\). This is
again an actual monomial line system, handled by HD-S34.

For \(A_r\), the root-line graph is the line graph of the complete
graph on \(r+1\) vertices. For \(r\ge4\), its maximal cliques
of size \(r\) are exactly the vertex stars; triangles have size
three. The argument of HD-S35 therefore gives every normalizer
as a scalar times an actual permutation of those vertices. The
Gaussian finite group containing \(W(A_r)\) and having character
field \(K\) must be the full \(\mu_4W(A_r)\). Apply the actual
scalar-augmentation factor and polarized maps of HD-S33.

The \(E_7\) and \(E_8\) diagrams have no nonidentity symmetry;
the three arms have pairwise distinct lengths. Their real normalizer
is therefore their Weyl group. The \(E_6\) diagram has precisely
one involution, interchanging its equal-length arms. In the Cartan
basis with edges
\((1,3),(3,4),(4,5),(5,6),(2,4)\), write
\(P=(1\ 6)(3\ 5)\) for that basis permutation. The exact root
certificate below gives an actual simple-reflection word \(w_0\)
with
\[
                            w_0=-P.                          \tag{S36.10}
\]
Thus its whole real normalizer is \(\{\pm w:w\in W(E_6)\}\).
Alternatively, \(\det C_6=3\) makes the discriminant group of
the root lattice cyclic of order three. Each root reflection acts
trivially on it, whereas \(-I\) does not. A Weyl element carrying
the opposite chamber to the chosen chamber consequently supplies
the nonidentity diagram symmetry after multiplication by \(-I\).
This gives the same explicit coset.

In all three exceptional real cases the Gaussian finite normalizer
is contained in \(\mu_4 W(E_r)\). Since it contains \(W(E_r)\),
each scalar component separately belongs to \(G'\); the exact
Gaussian character field forces \(iI\) and hence the full scalar
extension. HD-S27 gives the split determinant for \(E_6,E_7\),
and HD-S28 gives \([-1]\) or \([-3]\) for \(E_8\).
Every possibility in (S36.9) has now been connected to an actual
previously analyzed polarized source.

\paragraph{5. Rank at most eight and the precise remaining finite-group gate.}
For the original faithful rank-\(m\) sphere coefficient system,
\[
                       \sum_jc_j=2m+6.                       \tag{S36.11}
\]
Rank nine again requires \(b\ge12\). If \(G\) contains no
reflection, every nonidentity local monodromy has \(c_j\ge2\),
so \(2m+6\ge2b\ge24\). This is impossible for \(m\le8\).
Thus every source in those ranks is covered by the result already
proved. In any remaining finite Gaussian source with a different
determinant, the actual faithful coefficient group must be
reflection-free and have \(m\ge9\). At \(m=9\), equations
(S36.11) and \(b\ge12\) force exactly twelve inertia elements,
each of codimension two. This describes a strictly smaller remaining
source class; it does not assert its existence or its period rank.

The primitive-corner and fixed-polarization scope is the same throughout.
For a nonfaithful starting representation, quotient the original curve
by its kernel and use its pullback/norm maps as in HD-S35; the scalar
polarization multiplier has norm sixth power. Other primitive lines
use the actual cleared corner intertwiner. Compatible flat
polarization changes in a rank-nine family use HD-S17. No extra
Gaussian action on an unrelated special factor is included.

\paragraph{Whole-proof and finite checks.}
The complete proof was read by its author, including the equality of
idempotents (S36.5), low-genus automorphisms, the class-field inclusion,
and all normalizer maps. The exact script
\path{reflection_block_normalizer_certificate.py} reconstructs the
72,126,240 roots from the displayed \(E_6,E_7,E_8\) Cartan matrices,
checks their signs and diagram automorphisms, and gives the literal
36-reflection word for (S36.10). The chronological sequence of its
left-applied indices is
\[
\begin{gathered}
1,2,3,1,4,2,3,1,4,3,5,4,2,3,1,4,3,5,\\
4,2,6,5,4,2,3,1,4,3,5,4,2,6,5,4,3,1.
\end{gathered}
\]
The accompanying JSON saves the exact word matrix and Cartan matrix.
The universal block, covering, moduli, and polarized-factor arguments
are the proof above; they are not inferred from these finite checks.
